\silentchapter{Til lesaren}{til lesaren}
\begin{widebody}
\begin{foreordbody}

Boka har sju delar.

\textit{Ordlista} definerer dei tekniske termane i alfabetisk orden. Ho er til for å slå opp i undervegs, ikkje for å lesast fyrst.

\textit{Traktaten} er aksiomatisk. Ho byggjer eit omgrepsapparat frå grunnen av, gjennom nummererte proposisjonar i åtte seksjonar. Siste seksjon er den generative modellen, der grammatikk, landskap og lyskjegle blir sameinte i éi likning. Proposisjonane er merkte etter type: definisjon (\textit{d}), teorem (\textit{t}), postulat (\textit{a}), observasjon (\textit{o}), illustrasjon (\textit{i}).

\textit{Fundamentet} viser at rammeverket ikkje er isolert. Same formalismen er empirisk demonstrert over fjorten substratnivå, frå genregulatoriske nettverk til sivilisasjonar. Forskingsprogrammet som følgjer står skissert i fire punkt.

\textit{Etterordet} er materiellt og programmatisk i éin flyt. Frå larven og eingongskoppen via Palantir og merksemdsøkonomien til den praktiske krava til faget.

\textit{Tre brev} adresserer designaren, studenten og akademiet direkte.

\textit{Referansane} er nummererte. Superscript-tala i proposisjonane og etterordet peikar hit.

\textit{Notata} kommenterer enkeltproposisjonar, plasserer dei i tradisjonen, og listar falsifiseringsvilkåra.

Tre \textbf{kompositoriske spegelpar} held teksten saman. Rota i éin seksjon møter avslutninga i ein annan, på tvers av boka:

\begin{center}\footnotesize
\renewcommand{\arraystretch}{1.2}
\begin{tabular}{@{}p{28mm}p{24mm}p{28mm}@{}}
\toprule
\textbf{Proposisjon 1} & speglar & \textbf{7.3} \\
formverda er det som er tilfelle & & formverda er det som verkar \\[4pt]
\textbf{1.131} & speglar & \textbf{7.23} \\
ei form utan alternativ er nødvendigheit & & gapet mellom kapasitet og omhug er det predikerte sviktrommet \\[4pt]
\textbf{5.7} & speglar & \textbf{8.1} \\
kjerne-likninga, strukturell syntese & & generative modellen, probabilistisk syntese \\
\bottomrule
\end{tabular}
\end{center}

Akademikaren vil finne det meste i traktaten, fundamentet og notata. Studenten vil finne seg sjølv i etterordet. Designaren i praksis bør lese siste sida fyrst, og byrje på nytt.

\end{foreordbody}
\end{widebody}
